\section{Introduction}
\label{sec:intro}

EHR data is uneven: vitals go missing, values are physically impossible, records disagree across
joins, and observations go stale. Secondary use therefore begins with a data-quality-assurance pass
(Kahn et al., 2016), run by rule-checkers whose access is normally the entire record, and the habit
persists when the work is split across agents.

Least privilege is a design principle (Saltzer \& Schroeder, 1975). For covered entities it is a
duty under the HIPAA minimum necessary standard at 45 C.F.R. \S 164.514(d) (Privacy of Individually
Identifiable Health Information, 2013). That standard limits each class of worker to the protected
health information (PHI) its function requires. Here that means the eighteen Safe Harbor identifier
categories (U.S. Department of Health and Human Services, Office for Civil Rights, 2012).

The argument turns on four terms.

\begin{itemize}[leftmargin=1.4em, itemsep=0.15em]
  \item \textbf{Manifest.} The list of fields one agent may read.
  \item \textbf{Projection.} What enforces it, building a reduced copy of the record from only those
        fields, so a field left out never reaches the agent.
  \item \textbf{Width.} How much a manifest releases, swept here from minimal to full.
  \item \textbf{Consumed.} A field a check reads, as against one merely released to it.
\end{itemize}

Least privilege's enforcement mechanisms stop above the field (Section~\ref{sec:background}): a SMART
on FHIR scope (Mandel et al., 2016; HL7 International, 2024) reaches a resource type, never the
fields within one. No such scope can say that a plausibility check may see an observation's value but
not its subject reference. In the sources reviewed we found no pricing of field-level confinement on
a clinical task.

\textbf{Research question.} What does it cost to show each checker less of the record, and which
part of the system needs the patient's identity at all? \textbf{Aim.} To put a number on that cost,
on one data-quality task. \textbf{Significance.}
Least privilege below the resource type is at present asserted without a measured price, so a team
choosing between whole-record access and a manifest has no figure to weigh. This study supplies one,
on a single plausibility-centred task, three demonstration cohorts, one language model and three of
the four quality dimensions scored. It settles nothing outside that setting, and what it prices is
exposure, not re-identification risk.

Three contributions answer them. A field-granular privilege ablation prices what withholding one
consumed field costs, with a released-but-unconsumed field as its negative control: confining each
agent to a manifest costs nothing. But that zero is an identity rather than a measurement, since the
confined and unconfined systems read the same three fields. A tiered design cuts inter-agent exposure
by 98.6 per cent at no measured cost on the three dimensions scored. And the benchmark scores a
correct detection as a false positive, which bounds what the cost comparisons run on it can claim.
